% ChargeShield-FL --- DSN 2027 --- Introduction
% Portata in LaTeX 2026-09-16 (Sprint 10zz+115) dallo skeleton .docx, con
% tre correzioni sostanziali rispetto a quella versione:
%
% (1) Il .docx diceva che il sanity-check a cinque assi "rules out an
%     insensitive attack harness". Lo Sprint 10oo dichiarava esplicitamente
%     quel claim NON dimostrato, ed e' ora superato dal canary positive
%     control, che da' la risposta vera e piu' sfumata: lo strumento e'
%     sensibile sulla superficie raw-loss, la variante shadow-calibrata no.
% (2) Il .docx descriveva il Privacy Auditor come "real-time
%     membership-inference-risk auditor". Sprint 10zz+104 ha corretto questa
%     attribuzione: l'Auditor fa budget accounting e telemetria, NON leakage
%     detection, e non fa enforcement. Descrizione allineata.
% (3) Aggiunte le research question esplicite, che mancavano sia nel .docx
%     sia in LaTeX (quelle in docs/CaseStudies.md sono superate: chiedevano
%     "a quale epsilon l'attacco diventa indistinguibile dal caso", domanda
%     la cui premessa i risultati hanno smentito).

\section{Introduction}
\label{sec:introduction}

Federated learning (FL) is increasingly proposed as the privacy-preserving
path for training shared models over critical infrastructure --- smart
grids, EV charging networks, industrial control systems --- where
centralising raw operational data is both a regulatory liability and a
security risk. Whether FL delivers on that promise in realistic,
production-style deployments is an empirical question that the literature
answers piecemeal: one attack per paper, one dataset per paper, rarely a
production FL framework, rarely an industrial dataset.

The gap we address is narrower and, we argue, more consequential. Differential
privacy in FL is normally applied at \textbf{client-level}: clipping and noise
are attached to the update a client sends, because that update is the only
object crossing the communication channel. The attacks practitioners actually
fear, however, are \textbf{record-level}: they ask whether one particular
charging session was in someone's training set. A client-level
$(\varepsilon,\delta)$ guarantee does not bound record-level leakage when a
client holds thousands of records, and the size of that gap in practice is not
something theory settles. It has to be measured.

\subsection{Research Questions}
\label{sec:introduction-rq}

\begin{itemize}
\item[\textbf{RQ1}] \textbf{The mismatch.} When differential privacy is
calibrated at client level, how much protection does an individual record
actually receive against membership inference?

\item[\textbf{RQ2}] \textbf{The instrument.} When an attack reports no
leakage, is that a property of the system under test or of the attack? A
null privacy result is only informative if the measuring instrument is
independently known to be sensitive.

\item[\textbf{RQ3}] \textbf{The privilege.} How far does the answer depend
on where along the privatisation pipeline the adversary is assumed to
observe --- the raw update, the clipped update, or the noised one?
\end{itemize}

RQ2 is not methodological housekeeping. Most reported FL privacy nulls are
published without a positive control, which leaves them indistinguishable
from an attack that simply does not work; making that distinction
measurable is the part of this work we expect to transfer beyond EV
charging.

\subsection{Approach and Findings}

We present ChargeShield-FL, built around two domain-agnostic components: the
\emph{ML Plane}, an event-driven observability substrate that any FL training
loop can wire into without coupling the training code to whatever monitors it,
and the \emph{Privacy Auditor}, a config-driven consumer of those events that
tracks per-client privacy-budget consumption and raises operational alerts.
The Auditor is deliberately telemetry, not enforcement, and not a leakage
detector: it reports what the configured budget implies, while the question of
what actually leaks is answered by attacks, offline, in
Section~\ref{sec:framework-attacks}. Using real EV charging data (ACN-Data;
Caltech, JPL, Office~1) --- not synthetic data standing in for a deployment
--- we run three attacks across two independent observation surfaces, in
simulation and on a genuine multi-container NVFLARE/Containerlab federation.

On natural data we find no detectable membership signal in any configuration,
\emph{including with no DP at all}, and the null holds on all three scorers,
not only the shadow-calibrated one. Addressing RQ2, we then validate the
instrument instead of assuming it: injecting deliberately memorised canaries
recovers a large, reproducible signal on the raw-loss surface, which
establishes that the pipeline detects record-level memorisation when
memorisation exists. The same experiment shows that the per-record Gaussian
calibration of LiRA degenerates in this regime --- its variance estimates
collapse onto their floor --- so the scorer most often treated as the
strongest is, here, the least reliable.

Finally, we characterise \emph{why} the null arises, which turns a negative
result into a positive statement about when FL privacy audits can be expected
to find anything. The null is not explained by dataset size, and not by the
model being unable to memorise: it is a property of the representation. In our
six-feature session space, the median distance from a held-out session to its
nearest training session is $0.0146$, against $0.0149$ between training
sessions themselves --- a ratio of $0.98$. An adversary given the most
favourable purely geometric signal, the distance from the training set, would
achieve AUC $0.4935$. Members and non-members are mutually indistinguishable
\emph{before any model is trained}, which upper-bounds every attack
independently of the model, of DP, and of $\varepsilon$. Adding a
near-unique timestamp feature does not change this, because uniqueness is
symmetric across the two groups and therefore creates no asymmetry an attacker
can exploit.

\subsection{Contributions}

\begin{itemize}
\item The \emph{ML Plane} (Section~\ref{sec:framework-mlplane}): a
domain-agnostic, event-driven observability substrate for FL training that
decouples training code from the components observing it.

\item The \emph{Privacy Auditor} (Section~\ref{sec:framework-auditor}): a
config-driven budget-accounting and alerting consumer operating on any
node's update as a flat numeric dictionary, independent of domain and of
any specific attack.

\item An \emph{instrument-validation protocol} for FL privacy audits
(Section~\ref{sec:framework-canary}): a canary positive control with
per-seed random-initialisation baselines and a balanced role-exchange
negative control, together with the conditions under which each is valid
and the effective sample size each actually carries.

\item An empirical case study on real, multi-site EV charging data and a
containerised NVFLARE deployment, reporting a null that is accompanied by
the evidence needed to interpret it (Section~\ref{sec:results}).

\item A \emph{representation-level explanation} of the null, measured rather
than conjectured, that predicts when similar audits will and will not find
signal.
\end{itemize}
